% Part: first-order-logic
% Chapter: introduction
% Section: satisfaction

\documentclass[../../../include/open-logic-section]{subfiles}

\begin{document}

\olfileid{fol}{int}{sat}

\olsection{পরিতৃপ্তি}

!!{formula}s কী তা জানা হয়ে গেলে প্রথম-ক্রমের যুক্তিবিদ্যার অর্থতত্ত্বে
এখনই এগিয়ে যাওয়া যায়: এখানে মৌলিক সংজ্ঞাটি হলো !!a{structure}-এর।
আমাদের সরল ভাষার জন্য !!a{structure}~$\Struct M$-এর মাত্র তিনটি উপাদান:
\emph{!!{domain}} নামে একটি অশূন্য সেট~$\Domain{M}$,
$\Struct M$-এ $\Obj a$ যে বস্তুটিকে নির্দেশ করে এবং $\Struct M$-এ
$\Obj P$ যে বস্তুগুলির ক্ষেত্রে সত্য। $\Obj a$ যে বস্তুটিকে নির্দেশ করে
তাকে~$\Assign{\Obj a}{M}$ এবং $\Obj P$ যে বস্তুগুলির ক্ষেত্রে সত্য,
তাদের সেটকে~$\Assign{\Obj P}{M}$ দিয়ে প্রকাশ করা হয়। !!^a{structure}~$\Struct{M}$
ঠিক এই তিনটি বস্তু নিয়ে গঠিত: $\Domain{M}$,
$\Assign{\Obj a}{M} \in \Domain{M}$ এবং $\Assign{\Obj P}{M} \subseteq
\Domain{M}$। সাধারণ ক্ষেত্রটি আরও জটিল হবে, কারণ সেখানে অনেক
!!{predicate}s ও !!{constant}s থাকবে, !!{predicate}s-এর স্থানসংখ্যা একের
বেশি হতে পারে এবং !!{function}s-ও থাকবে।% BN-SRC-085 TARGET-CORRECTION-BEGIN
\par\smallskip\noindent\textnormal{\small\textbf{উৎস-সংশোধন (BN-SRC-085):} হিমায়িত ইংরেজি উৎসে সাধারণ ক্ষেত্রটির বর্ণনায় বহু বিধেয় ও ধ্রুবকের পরে ভুল করে ধ্রুবকগুলির স্থানসংখ্যা একের বেশি হতে পারে বলা হয়েছে। সম্পর্কীয় স্থানসংখ্যা বিধেয়ের ধর্ম; ধ্রুবক শূন্য-স্থানীয়। বাংলা পাঠে তাই বিধেয়গুলির স্থানসংখ্যাই একের বেশি হতে পারে বলা হয়েছে।} % BN-SRC-085 TARGET-CORRECTION-END

যেসব !!{formula}s-এ !!{variable}s নেই, তাদের পরিতৃপ্তির সংজ্ঞা দেওয়ার
জন্য এটুকুই যথেষ্ট। আমরা যে পথে !!{formula}s সংজ্ঞায়িত করেছি, তার
প্রতিচ্ছবি এমন একটি আরোহী সংজ্ঞা দেওয়াই মূল ধারণা। কোনো পরমাণু সূত্র
কখন~$\Struct{M}$-এ পরিতৃপ্ত হয় তা নির্দিষ্ট করি; তারপর, যেমন, $!A$
$\Struct{M}$-এ পরিতৃপ্ত হয় কি না, তার ভিত্তিতে $\lnot !A$ কখন
$\Struct{M}$-এ পরিতৃপ্ত হয় তা বলি। যেমন, আমরা এভাবে সংজ্ঞায়িত করতে
পারি:
\begin{enumerate}
  \item $\Atom{\Obj P}{\Obj a}$,~$\Struct{M}$-এ পরিতৃপ্ত হয় যদি এবং
  কেবল যদি $\Assign{\Obj a}{M} \in \Assign{\Obj P}{M}$।
  \item $\lnot !A$,~$\Struct{M}$-এ পরিতৃপ্ত হয় যদি এবং কেবল যদি $!A$,
  $\Struct{M}$-এ পরিতৃপ্ত না হয়।
  \item $(!A \land !B)$,~$\Struct{M}$-এ পরিতৃপ্ত হয় যদি এবং কেবল যদি
  $!A$,~$\Struct{M}$-এ পরিতৃপ্ত হয় এবং $!B$-ও~$\Struct{M}$-এ পরিতৃপ্ত
  হয়।
\end{enumerate}
ধরা যাক $\Domain{M} = \{0, 1, 2\}$, $\Assign{\Obj a}{M} = 1$ এবং
$\Assign{\Obj P}{M} = \{1, 2\}$। এই সংজ্ঞা থেকে পাই,
$\Atom{\Obj P}{\Obj a}$,~$\Struct{M}$-এ পরিতৃপ্ত হয় (কারণ
$\Assign{\Obj a}{M} =  1 \in \{1,2\} = \Assign{\Obj P}{M}$)। আরও পাই,
$\lnot \Atom{\Obj P}{\Obj a}$,~$\Struct{M}$-এ পরিতৃপ্ত হয় না; আবার তার
ফলে $\lnot\lnot \Atom{\Obj P}{\Obj a}$ পরিতৃপ্ত হয় এবং
$(\lnot \Atom{\Obj P}{\Obj a} \land \Atom{\Obj P}{\Obj a})$ পরিতৃপ্ত
হয় না; এভাবে চলতে থাকে।

পরিমাণসূচকের সংজ্ঞা দিতে গেলেই অসুবিধা দেখা দেয়: আমরা এমন কিছু বলতে চাই,
“$\lexists[\Obj v_0][\Atom{\Obj P}{\Obj v_0}]$ পরিতৃপ্ত হয় যদি এবং
কেবল যদি $\Atom{\Obj P}{\Obj v_0}$ পরিতৃপ্ত হয়।” কিন্তু
!!{structure}~$\Struct{M}$, !!{variable}s নিয়ে কী করতে হবে তা বলে না।
আমরা আসলে বলতে চাই, $\Obj v_0$-এর \emph{কোনো একটি মানের জন্য}
$\Atom{\Obj P}{\Obj v_0}$ পরিতৃপ্ত হয়। একে নির্ভুল করতে শুধু $\Obj a$-কে
নয়, $\Obj v_0$-কেও~$\Domain{M}$-এর !!{element}s দেওয়ার একটি উপায়
দরকার। সেই উদ্দেশ্যে আমরা !!{variable}-এর \emph{আরোপ} প্রবর্তন করি।
!!^a{variable} আরোপ কেবল এমন একটি অপেক্ষক~$s$, যা !!{variable}s-কে
$\Domain{M}$-এর !!{element}s-এ চিত্রিত করে (আমাদের উদাহরণে $0$, $1$ বা~$2$-এর
কোনো একটিতে)। কোন কোন !!{variable}s কোনো !!a{formula}-তে দেখা দিতে পারে
তা আগে থেকে জানা নেই, তাই $s$ কোন !!{variable}s-কে মান দেবে তা সীমিত
করা যায় না। সহজ সমাধান হলো, $s$ যেন \emph{সব} !!{variable}s~$\Obj v_0$,
$\Obj v_1$, \dots-কে মান দেয়—এই শর্ত রাখা। আমরা শুধু প্রয়োজনীয়গুলিই
ব্যবহার করব।% BN-SRC-086 TARGET-CORRECTION-BEGIN
\par\smallskip\noindent\textnormal{\small\textbf{উৎস-সংশোধন (BN-SRC-086):} হিমায়িত ইংরেজি উৎসে $\Domain{M}=\{0,1,2\}$ স্থির করার পরও চলরাশি-আরোপের সম্ভাব্য মানকে $1$, $2$ বা $3$ বলা হয়েছে। বাংলা পাঠে ঘোষিত সংজ্ঞাক্ষেত্রের ঠিক তিন সদস্য $0$, $1$ ও $2$ রাখা হয়েছে।} % BN-SRC-086 TARGET-CORRECTION-END

!!{formula}s-এর পরিতৃপ্তিকে শুধু !!a{structure}-এর সাপেক্ষে সংজ্ঞায়িত না
করে, !!a{structure}~$\Struct{M}$ \emph{ও} !!a{variable} আরোপ~$s$—দুটির
সাপেক্ষে সংজ্ঞায়িত করব; সংক্ষেপে লিখব $\Sat{M}{!A}[s]$। এখন আমাদের
সংজ্ঞায় !!{variable}s-সহ পরমাণু !!{formula}s সামলানোর একটি অতিরিক্ত
অনুচ্ছেদ থাকবে:
\begin{enumerate}
  \item $\Sat{M}{\Atom{\Obj P}{\Obj a}}[s]$ যদি এবং কেবল যদি
  $\Assign{\Obj a}{M} \in \Assign{\Obj P}{M}$।
  \item $\Sat{M}{\Atom{\Obj P}{\Obj v_i}}[s]$ যদি এবং কেবল যদি
  $s(\Obj v_i) \in \Assign{\Obj P}{M}$।
  \item $\Sat{M}{\lnot !A}[s]$ যদি এবং কেবল যদি $\Sat{M}{!A}[s]$ নয়।
  \item $\Sat{M}{(!A \land !B)}[s]$ যদি এবং কেবল যদি $\Sat{M}{!A}[s]$
  এবং $\Sat{M}{!B}[s]$।
\end{enumerate}
বেশ, এতে একটি সমস্যার সমাধান হলো: $s(\Obj v_0)$ মানটির জন্য
$\Struct{M}$ কখন $\Atom{\Obj P}{\Obj v_0}$-কে পরিতৃপ্ত করে, তা এখন বলা
যায়। $\lexists[\Obj v_0][\Atom{\Obj P}{\Obj v_0}]$-এর সংজ্ঞাটি ঠিক
করতে আরও একটি কাজ করতে হবে। আমরা চাই,
$\Sat{M}{\lexists[\Obj v_0][\Atom{\Obj P}{\Obj v_0}]}[s]$ হবে যদি এবং
কেবল যদি $\Obj v_0$-কে কোনো-একভাবে মান দেয় এমন \emph{কোনো} $s'$-এর
জন্য $\Sat{M}{\Atom{\Obj P}{\Obj v_0}}[s']$ হয়। কিন্তু $\Obj v_0$-কে
দেওয়া সেই মানটি $s(\Obj v_0)$-এর নির্দেশিত মান হতেই হবে না। এর জন্য
একটি সংকেত প্রবর্তন করব: $m \in \Domain{M}$ হলে
$\Subst{s}{m}{\Obj v_0}$ হবে এমন আরোপ, যা~$s$-এরই মতো (অর্থাৎ~$\Obj v_0$
ছাড়া অন্য সব !!{variable}s-এর জন্য একই), শুধু $\Obj v_0$-কে সেটি~$m$
মান দেয়। এখন আমাদের সংজ্ঞাটি হতে পারে:
\begin{enumerate}\setcounter{enumi}{4}
  \item $\Sat{M}{\lexists[\Obj v_i][!A]}[s]$ যদি এবং কেবল যদি কোনো
  $m \in \Domain{M}$-এর জন্য $\Sat{M}{!A}[\Subst{s}{m}{\Obj v_i}]$।
\end{enumerate}
সংজ্ঞাটি কি ঠিকমতো কাজ করে? ধরা যাক, সব~$i \in \Nat$-এর জন্য
$s(\Obj v_i) = 0$। $\Sat{M}{\lexists[\Obj v_0][\Atom{\Obj P}{\Obj
v_0}]}[s]$ হবে যদি এবং কেবল যদি এমন $m \in \Domain{M}$ থাকে যাতে
$\Sat{M}{\Atom{\Obj P}{\Obj v_0}}[\Subst{s}{m}{\Obj v_0}]$। এমন মান
আছে: $m = 1$ বা $m = 2$ বেছে নেওয়া যায়। লক্ষ করো, $s$ নিজে
$\Obj v_0$-কে যে মান~$s(\Obj v_0)$ দেয়—এক্ষেত্রে $0$—তা কাজটি না করলেও
বাক্যটি সত্য। আমাদের $\Sat{M}{\Atom{\Obj P}{\Obj v_0}}[\Subst{s}{1}{\Obj
v_0}]$ আছে, কিন্তু $\Sat{M}{\Atom{\Obj P}{\Obj v_0}}[s]$ নেই।

এটি যদি বিভ্রান্তিকর ও জটিল মনে হয়, তার কারণ এটি সত্যিই তেমন। কিন্তু সব
!!{formula}s-এর জন্য পরিতৃপ্তির একটি নির্ভুল আরোহী সংজ্ঞা দিতে এই বাড়তি
জটিলতা দরকার, এবং অর্থগত ধারণাগুলিকে নির্ভুলভাবে সংজ্ঞায়িত করতেও এমন
কিছু প্রয়োজন। কাজটি অন্যভাবেও করা যায়, তবে প্রতিটি উপায়ই সমানভাবে
অসুন্দর কিংবা সুন্দর।

\end{document}
